\documentclass[10pt,a4paper]{extarticle}
\pagenumbering{gobble}
\usepackage[utf8]{inputenc}
\usepackage[T1]{fontenc}
\usepackage[left=0.6in,top=0.6in,right=0.6in,bottom=0.6in]{geometry}
\usepackage{hyperref}
\usepackage{titlesec}
\usepackage{enumitem}
\usepackage{graphicx}
\titleformat{\section}{\large\bfseries\scshape\raggedright}{}{0em}{}[\titlerule]
\titlespacing*{\section}{0pt}{12pt}{8pt}
\begin{document}
\begin{center}
    \begin{minipage}{\textwidth}
        \centering
        {\LARGE\textbf{Emre Kolunsağ}} \hspace{2pt} {\LARGE{Jr. SAP ABAP Danışmanı / Yazılım Mühendisi}}\\[10pt]
        İstanbul, Türkiye \textbullet\
        \href{mailto:emrekolunsag@gmail.com}{emrekolunsag@gmail.com} \textbullet\
        \href{https://linkedin.com/in/emrekolunsag}{linkedin.com/in/emrekolunsag} \textbullet\
        \href{https://github.com/kolsagg}{github.com/kolsagg}
    \end{minipage}
\end{center}
\section{Özet}
Yazılım Mühendisliği mezunu, 1 yılı aşkın süredir SAP ABAP Danışmanı olarak kurumsal projelerde görev alan
deneyimli geliştirici. SAP S/4HANA Greenfield ve migration (geçiş) projelerinde; BTE, User-Exit, Enhancement,
BAPI ve Web Services (REST/SOAP) teknolojileri ile SAP uyumlu teknik çözümlerin ve modül geliştirmelerinin
tasarlanması, optimize edilmesi ve devreye alınması konusunda uzmanlaşmıştır. Sanal POS ve e-Dönüşüm
entegrasyonlarıyla FI modülü süreçlerinin otomatize edilmesi konusunda tecrübe sahibidir. Modern web/yazılım
geliştirme (TypeScript, Python, SQL) altyapısı sayesinde SAP ve dış sistem entegrasyonlarına bütüncül ve
proaktif bir problem çözme yaklaşımı getirmektedir.
\section{Yetenekler}
\textbf{SAP \& ABAP Teknolojileri:} ABAP Programming, ABAP Objects, Open SQL, BTE, User-Exit, Enhancements, BAPI, Web Services (REST/SOAP), SAP FI Modülü\\
\textbf{SAP Geçiş \& Entegrasyon:} SAP S/4HANA Greenfield Implementation, S/4HANA Migration, Sanal POS Entegrasyonu, Elektronik Çözüm Entegrasyonu (e-Dönüşüm), Dönemselleştirme (Accrual/Deferral)\\
\textbf{Geliştirme Araçları \& Veritabanı:} SAP GUI, HANA DB, PostgreSQL, MySQL, Supabase, Git, Postman, SOAPUI, REST API\\
\textbf{Diğer Programlama Dilleri:} TypeScript, Python, JavaScript, Java, HTML/CSS, React.js, Node.js\\
\textbf{Yabancı Dil:} Türkçe (Ana dili), İngilizce (Orta Seviye)
\section{İş Deneyimi}
\textbf{ConcentIT. \textit{(İstanbul, Türkiye)}}\\
\textit{Jr. SAP ABAP Danışmanı / Geliştirici} \hfill Aralık 2024 -- Halen
\begin{itemize}[leftmargin=*,noitemsep,topsep=0pt]
    \item SAP S/4HANA geçiş ve greenfield uyarlama projelerinde müşteri gereksinimlerine uygun özel ABAP programları, raporlar ve iyileştirmeler (enhancements) geliştirildi.
    \item Müşteri süreçlerindeki sistem kesintilerini en aza indirmek için ABAP kodlarında hata ayıklama (debugging), performans optimizasyonu ve teknik destek süreçleri yürütüldü.
    \item FI modülü başta olmak üzere iş süreçlerini destekleyen custom entegrasyon ve otomasyon çözümleri tasarlandı.
\end{itemize}
\textbf{Kartelam, Fabriqa Inovation A.Ş. \textit{(İstanbul, Türkiye)}}\\
\textit{Uzun Dönem Stajyer -- Full-Stack Geliştirici} \hfill Ekim 2023 -- Kasım 2024
\begin{itemize}[leftmargin=*,noitemsep,topsep=0pt]
    \item Ölçeklenebilir web uygulamaları ve backend mimarileri geliştirilerek yazılım mühendisliği prensipleri ve kod kalitesi standartları uygulandı.
    \item Farklı disiplinlerden ekiplerle çalışılarak canlı ortamlara dağıtım (deployment) ve sistem entegrasyonu süreçleri desteklendi.
\end{itemize}
\section{SAP Proje Deneyimleri}
\textbf{Coşkunöz Holding \textnormal{--- \textit{S/4HANA Greenfield İmplementasyonu}}}
\begin{itemize}[leftmargin=*,noitemsep,topsep=0pt]
    \item Uçtan uca iş süreçlerine yönelik özel ABAP programları ve ALV raporları geliştirildi
\end{itemize}
\textbf{İlbak Holding \textnormal{--- \textit{S/4HANA Geçişi (Migration)}}}
\begin{itemize}[leftmargin=*,noitemsep,topsep=0pt]
    \item Custom program ve formlar yeni veri modeline taşındı
    \item S/4HANA uyumluluk düzeltmeleri ve dönüşüm sonrası teknik testler yürütüldü
\end{itemize}
\textbf{Limak Holding \textnormal{--- \textit{Sanal POS Entegrasyonu (FI)}}}
\begin{itemize}[leftmargin=*,noitemsep,topsep=0pt]
    \item Sanal POS ödeme portalına açık faturaları ileten REST Web Service geliştirildi
    \item Ödemeleri otomatik muhasebeleştiren ABAP Background Job (arka plan işi) programı geliştirildi
    \item Sürecin kolay izlenebilmesi için FI modülünde ALV rapor ve kokpit ekranı geliştirildi
\end{itemize}
\textbf{Kılıç Deniz \textnormal{--- \textit{Elektronik Çözüm (e-Dönüşüm) Entegrasyonu}}}
\begin{itemize}[leftmargin=*,noitemsep,topsep=0pt]
    \item Harici sistem -- SAP veri akışı için Web Service ve BTE/User-Exit geliştirildi
    \item FI modülü için özel program ve form uyarlamaları yapıldı
\end{itemize}
\textbf{JCI \textnormal{--- \textit{Dönemselleştirme Süreçleri}}}
\begin{itemize}[leftmargin=*,noitemsep,topsep=0pt]
    \item Finansal hesaplamaları otomatize eden ABAP çözümleri geliştirildi
    \item Dönemselleştirme mantığına uygun iş süreci iyileştirmeleri yapıldı
\end{itemize}
\textbf{Bilkent Holding \textnormal{--- \textit{Destek Hizmetleri}}}
\begin{itemize}[leftmargin=*,noitemsep,topsep=0pt]
    \item FI modülü için ABAP teknik desteği sağlandı
    \item Mevcut program/raporlar geliştirildi, SmartForms ve enhancement düzenlemeleri yapıldı
\end{itemize}
\textbf{Limak Holding \textnormal{--- \textit{Destek Hizmetleri}}}
\begin{itemize}[leftmargin=*,noitemsep,topsep=0pt]
    \item Rutin ABAP geliştirme ve debugging tabanlı hata giderme yapıldı
    \item İş süreci optimizasyonu desteği verildi
\end{itemize}
\textbf{Atakaş Holding \textnormal{--- \textit{Destek Hizmetleri}}}
\begin{itemize}[leftmargin=*,noitemsep,topsep=0pt]
    \item Sürekli ABAP teknik desteği sağlandı
    \item Canlı sistemde rapor, program ve exit geliştirmeleri yürütüldü
\end{itemize}
\textbf{AGT \textnormal{--- \textit{Destek Hizmetleri}}}
\begin{itemize}[leftmargin=*,noitemsep,topsep=0pt]
    \item Ağırlıklı olarak FI modülü kapsamında ABAP geliştirme desteği verildi
    \item Form/rapor iyileştirmeleri ve iş süreci optimizasyonları yapıldı
\end{itemize}
\textbf{IATCO \textnormal{--- \textit{Destek Hizmetleri}}}
\begin{itemize}[leftmargin=*,noitemsep,topsep=0pt]
    \item Sürekli ABAP desteği ve FI modülü geliştirmeleri sağlandı
    \item Canlı sistemdeki teknik sorunlar debugging ile çözüldü
\end{itemize}
\section{Diğer Yazılım Projeleri}
\textbf{Yapay Zeka Destekli Finans Takip Mobil Uygulaması} \hfill \textit{Flutter, Python, Supabase, Ollama}
\begin{itemize}[leftmargin=*,noitemsep,topsep=0pt]
    \item QR fiş tarama ve yapay zeka destekli öneriler sunan harcama takip mobil uygulaması geliştirildi
    \item Python ile JSON formatındaki kullanıcı verilerini AES-256 ile şifreleyen güvenli backend tasarlandı
    \item Flutter ile responsive ve kullanıcı dostu arayüz geliştirildi
    \item Kimlik doğrulama ve veritabanı yönetimi için Supabase entegre edildi
    \item Proje geliştirme sürecinde Agile (Scrum) metodolojisi uygulandı
\end{itemize}
\section{Eğitim}
\textbf{İstanbul Aydın Üniversitesi} \hfill İstanbul, Türkiye\\
\textit{Lisans, Yazılım Mühendisliği (İngilizce)} \hfill 2020 -- 2025 (Mezuniyet: Ekim 2025)
\end{document}
